\section{Mother’s Day 2020}

\begin{quotex}
The truth is that memory does not consist in a regression from the present to the past, but on the contrary in a progress from the past to the present. ... pure memory is a spiritual manifestation. With memory we are in very truth in the domain of the spirit, the state of the brain continues the remembrance; it gives it a hold on the present by the materiality which it (acting as a mirror) confers upon it.
\flright{\textsc{Henri Bergson}, \emph{Matter and Memory}}
\end{quotex}


Perception provides an immediate experience of the world, but memory creates a world through the power of memory to retain and recall past perceptions and images. Memory is therefore actual in the present and does not mean a return to the past. A good spiritual exercise is the attempt to remember because the one's past needs to be integrated into, and made present to, the Self.

The inability to remember is a curse. In her final years, my Mother was afflicted with Alzheimer's, which frustrated and confused her. She had the intelligence and wherewithal to be able to pretend. My father posted picture of me and my children on the wall and explained them to her just prior to my visits. She still had a vague memory of me, but her acting ability, when asking about my kids, was impressive.

\paragraph{The Anglophile}


My Mother was an anglophile. The most visible symbol of that is that she never made Sunday Gravy\footnote{\url{https://www.afamilyfeast.com/sunday-gravy/}}, as did the rest of the family. Instead, that day was reserved for the Sunday Roast\footnote{\url{https://en.wikipedia.org/wiki/Sunday_roast}}, which is a British tradition. The result is that we developed a broader, less clannish, perspective on the world.

\paragraph{Raising the Scientist}


Hence, I have often looked back to understand how she raised me. The disconnected images and memories have gradually coalesced into a coherent view. My uncle --- her brother --- was in the field of education and would give us (also my sister) IQ tests. Since my father was a chemist, they must have decided to prep me for one of the STEM fields\footnote{\url{https://www.britannica.com/topic/STEM-education}}, years before it was a ``thing''.

I recall that she would give me long strings of numbers to add up, even before I started school. I always had models to build --- jets, missiles, rocket ships, automobiles ---. Toy stores had aisle of scale models until the alleged glue sniffing craze destroyed that industry.

There were no malls, so a visit to the department store was not a common event. I recall that she would always bring me a new Tom Corbett, Space Cadet\footnote{\url{https://en.wikipedia.org/wiki/Tom_Corbett,_Space_Cadet}} novel. I had subscriptions to I Spy\footnote{\url{https://en.wikipedia.org/wiki/I-Spy_(Michelin)}} and Science Digest from an early age. Of course, there were books about Dinosaurs, Outer Space, etc., which many boys still like today.

When I was a little older, around 10, I went to classes at the Boston Museum of Science\footnote{\url{https://www.mos.org/}} on Saturday mornings as they shopped in town. After class, I would wait in the library until they arrived. My favorite books were about Astronomy and Mathematics. The idea that those strange mathematical symbols could encapsulate so much information excited and intrigued me.

Rene Descartes fascinated me. I recall reading that, as a boy, he was allowed to remain in bed in the mornings because of his intelligence. Although I did not get that particular indulgence, I aspired to me a mathematician-philosopher.

\paragraph{The Poison of Liberalism}


My sister did not get the same training despite her high IQ. She was a pretty girl and my mother used to sew nice dresses for her. Instead of Tom Corbett, she got Honey Bunch books.

One day, I related this story to someone who knew both my mother and me. Instead of seeing a loving mother who wanted the best life for her son, this person's immediate reaction was that my mother was a ``sexist''. That truly upset me, and I seldom get upset, so I called her out. Unfortunately, ideologues like that simply lack the ability to see the ``person'' that my mother was. Instead, she chose to objectify her ss one of the few Ready Made Ideas\footnote{\url{https://gornahoor.net/?p=11832}} that form the substance of her distorted worldview. As obstinate as she was, she refused to back down. Ultimately I was the one who was forced to apologize in order to keep the peace.


\flrightit{Posted on 2020-05-10 by Cologero}

